% =====================================================================
% TRACE-ESUS — formal specification of the Experiment 1 generator
% Paste target: bioRxiv Paper 1, Methods (§4 "Synthetic data-generating
% process"). Matches r21_latent_endotyping_experiment.py exactly:
%   simulate_two_mechanism_cohort()  (line 232)
%   SimulationConfig                 (line 77)
% Requires: amsmath, amssymb, bm
% =====================================================================

\subsection{Structural causal model}

We specify the data-generating process as a structural causal model
$\mathcal{M} = \langle \mathbf{U}, \mathbf{V}, \mathcal{F}, P(\mathbf{U}) \rangle$
with endogenous variables $\mathbf{V} = \{R, Z, \mathbf{B}\}$ and mutually
independent exogenous variables $\mathbf{U} = \{U_R, U_Z, \bm{\varepsilon}\}$.
Here $R \in \{0,1\}$ denotes renal dysfunction, $Z \in \{a, c\}$ the latent
mechanism (atrial-thromboembolic vs.\ competing), and
$\mathbf{B} = (B_1, B_2, B_3)^\top \in \mathbb{R}^3$ the biomarker vector
(NT-proBNP-like, atrial electrical evidence, competing-mechanism evidence).
The structural equations $\mathcal{F}$ are

\begin{align}
  R &:= \mathbb{1}\{U_R < \pi_R\},
    && U_R \sim \mathrm{Unif}(0,1), \label{eq:scm-r}\\
  Z &:= a \cdot \mathbb{1}\{U_Z < p_R\} + c \cdot \mathbb{1}\{U_Z \geq p_R\},
    && U_Z \sim \mathrm{Unif}(0,1), \label{eq:scm-z}\\
  \mathbf{B} &:= \bm{\mu}_Z + R\,\bm{\delta} + \bm{\varepsilon},
    && \bm{\varepsilon} \sim \mathcal{N}(\mathbf{0}, \bm{\Sigma}),
    \label{eq:scm-b}
\end{align}

where the assignment symbol $:=$ denotes a structural (asymmetric)
relation rather than an algebraic equality. All effects are expressed in
units of within-biomarker residual standard deviation:

\begin{equation}
  \bm{\mu}_a = \begin{pmatrix} 1.25 \\ 1.00 \\ 0 \end{pmatrix}, \quad
  \bm{\mu}_c = \begin{pmatrix} 0 \\ 0 \\ 1.00 \end{pmatrix}, \quad
  \bm{\delta} = \begin{pmatrix} \delta_R \\ 0 \\ 0 \end{pmatrix}, \quad
  \bm{\Sigma} = \mathbf{I}_3 ,
  \label{eq:params}
\end{equation}

with renal prevalence $\pi_R = 0.30$, mechanism prevalence
$p_0 = p_1 = 1/2$, and renal distortion strength
$\delta_R \in \{0,\ 0.5,\ 1.0,\ 1.5\}$ indexing the four confounding
levels (None, Weak, Moderate, Strong).

The induced directed acyclic graph is
$\mathcal{G} = (\mathbf{V}, \mathcal{E})$ with edge set
\begin{equation}
  \mathcal{E} = \{ Z \rightarrow B_j \}_{j=1,2,3}
              \ \cup\ \{ R \rightarrow B_1 \}.
  \label{eq:edges}
\end{equation}
The edge $R \rightarrow Z$ is \emph{absent by construction}
($p_0 = p_1$). Renal dysfunction is therefore not a confounder of the
mechanism--biomarker relation in the classical sense; it is a
\emph{competing cause of a shared measurement}, colliding with the
mechanism path at $B_1$. This distinction is deliberate and is what the
modular formulation is designed to address.

\subsection{Markov factorization and the observed mixture}

The joint distribution factorizes according to $\mathcal{G}$ as
\begin{equation}
  P(R, Z, \mathbf{B})
    = P(R)\, P(Z \mid R)\, P(\mathbf{B} \mid Z, R),
  \label{eq:factorization}
\end{equation}
with $P(Z \mid R) = P(Z)$ under \eqref{eq:params}. Marginalizing over
both $R$ and $Z$ yields the distribution actually available to an
unsupervised learner:
\begin{equation}
  P(\mathbf{B}) = \sum_{r \in \{0,1\}} \sum_{z \in \{a,c\}}
    P(R{=}r)\, p_r(z)\;
    \mathcal{N}\!\left( \mathbf{B};\ \bm{\mu}_z + r\bm{\delta},\ \bm{\Sigma} \right).
  \label{eq:four-component}
\end{equation}
Equation~\eqref{eq:four-component} is a \emph{four}-component Gaussian
mixture: two mechanisms crossed with two renal states. Only the
mechanism partition is of scientific interest. A two-class mixture model
fitted to \eqref{eq:four-component} without a renal path must select one
of these partitions, and for sufficiently large $\delta_R$ the renal
partition attains the higher likelihood. This is the formal basis of the
false-endotype-discovery result reported in Section~[K1-NULL].

\subsection{Bayes-optimal mechanism posterior}

Because the two components share $\bm{\Sigma}$, the posterior under the
data-generating model has closed linear-discriminant form. Writing
$\mathbf{w} = \bm{\Sigma}^{-1}(\bm{\mu}_a - \bm{\mu}_c)$ and
$\bar{\bm{\mu}} = \tfrac{1}{2}(\bm{\mu}_a + \bm{\mu}_c)$,
\begin{equation}
  P(Z = a \mid \mathbf{B}, R)
    = \sigma\!\left(
        \operatorname{logit} p_R
        + \mathbf{w}^\top \left( \mathbf{B} - R\bm{\delta} - \bar{\bm{\mu}} \right)
      \right),
  \label{eq:bayes-posterior}
\end{equation}
where $\sigma(\cdot)$ is the logistic function. Equation
\eqref{eq:bayes-posterior} makes the role of the nuisance path explicit:
correct inference requires only that the term $R\bm{\delta}$ be removed
before the discriminant is applied. Under \eqref{eq:params} the
Mahalanobis separation is
\begin{equation}
  d = \left\| \bm{\mu}_a - \bm{\mu}_c \right\|_{\bm{\Sigma}^{-1}}
    = \sqrt{1.25^2 + 1^2 + 1^2} = 1.8875,
  \label{eq:separation}
\end{equation}
so with $p_R = 1/2$ the Bayes error is
$\Phi(-d/2) = 0.1727$ and the attainable accuracy ceiling is
$0.8273$, invariant to $\delta_R$. Simulated oracle performance matches
this analytic value to four decimal places at every distortion level,
confirming that the renal path is fully removable given $R$.

\subsection{Model family and the causal constraint}

All fitted models are two-component mixtures differing only in their
treatment of the nuisance path. Let $\mathcal{M} \in \{0,1\}^3$ denote a
prespecified support mask on $\bm{\delta}$. We compare:

\begin{enumerate}
  \item \textbf{Pooled associative latent class model:}
        $P(Z)\,P(R \mid Z)\,P(\mathbf{B} \mid Z)$, i.e.\ $\bm{\delta} \equiv \mathbf{0}$
        with $R$ modeled as a \emph{child} of the latent class.
  \item \textbf{Renal-adjusted associative model:}
        $P(Z \mid R)\,P(\mathbf{B} \mid Z, R)$ with $\mathcal{M} = (1,1,1)$
        (all nuisance paths free).
  \item \textbf{Biology-constrained latent SCM:}
        $P(Z \mid R)\,P(\mathbf{B} \mid Z, R)$ with $\mathcal{M} = (1,0,0)$,
        the true support, fixed a priori from biology rather than estimated.
\end{enumerate}

The causal model is thus a \emph{prespecified sparsity constraint on the
nuisance-path support}. A data-generating oracle with known
$(\bm{\mu}_a, \bm{\mu}_c, \bm{\delta}, \bm{\Sigma}, p_R)$ provides the
performance ceiling \eqref{eq:separation}.

\subsection{Testable conditional independences}

The graph \eqref{eq:edges} implies three statements, of which the third
is directly checkable in observational data and serves as a positive
control for the anchoring procedure:
\begin{align}
  Z &\perp\!\!\!\perp R && \text{(no edge; no common cause)} \label{eq:dsep1}\\
  Z &\not\perp\!\!\!\perp R \mid B_1 && \text{(collider at $B_1$)} \label{eq:dsep2}\\
  (B_2, B_3) &\perp\!\!\!\perp R \mid Z && \text{(no renal path to the anchors)} \label{eq:dsep3}
\end{align}
Statement \eqref{eq:dsep2} is the formal explanation for the failure of
the pooled model: conditioning on the biomarkers, as any mixture model
does by construction, induces dependence between mechanism and renal
status at the collider, which the term $P(R \mid Z)$ then absorbs and
subsequently uses as evidence. Statement \eqref{eq:dsep3} justifies
orienting latent labels by the contrast $B_2 - B_3$, which excludes the
renal-contaminated marker $B_1$.

\subsection{Counterfactual quantities}

Because $\bm{\varepsilon}$ enters \eqref{eq:scm-b} additively and is
shared across mechanism branches, the potential outcomes are
\begin{equation}
  \mathbf{B}(z) := \bm{\mu}_z + R\bm{\delta} + \bm{\varepsilon},
  \qquad
  \mathbf{B}(z, r) := \bm{\mu}_z + r\bm{\delta} + \bm{\varepsilon}.
  \label{eq:potential-outcomes}
\end{equation}
Abduction recovers
$\bm{\varepsilon} = \mathbf{B} - \bm{\mu}_z - R\bm{\delta}$, which depends
on the latent $z$; counterfactual scores are therefore computed
branch-wise and averaged under $P(Z \mid \mathbf{B}, R)$ rather than by
abducting a single residual under an assumed branch.

\subsection{Assumptions encoded in the generator}

The specification \eqref{eq:scm-r}--\eqref{eq:params} deliberately
assumes: (i) \emph{additivity} --- the renal contribution is invariant to
mechanism, with no $Z \times R$ interaction; (ii) \emph{homoscedasticity}
--- $\bm{\Sigma}$ does not depend on $Z$ or $R$; (iii) \emph{mutually
exclusive mechanisms} --- $Z$ is scalar-valued, so no patient carries
both pathologies; (iv) $K = 2$; (v) complete data, with no missingness or
assay heterogeneity (these are introduced in the transportability
experiment). Assumption (iii) is consequential: with mutually exclusive
mechanisms no alternative causal path survives disablement, so the
redundancy phenomenon the counterfactual framework is designed to detect
cannot arise under this generator. Relaxing (i) and (iii) is the subject
of Section~[DIVERGENCE].
